\section{Growth of the differential and decay of the exact part}\label{sec:decay}

In this section $\mathsf E$ is a complex with a surjective differential at
$i_0$ and $r\geq1$. The estimates of Lemma~\ref{lem:estimates} and the proof
of Theorem~\ref{thm:decay} follow those of \cite{wue26b} for $K_m$; what is
new is that the block $z_\ast+\mu_{i_0}$ is no longer $z_\ast$ plus a
positive matrix, and that the phases at the targets of the arrows out of
$i_0$ are bounded below only through the a priori bound. Recall $\lambda=\lambda_{\min}(L)=\lambda_{\min}(F^\dagger F)$, and
let
\[
 \varpi_\star=\min\bigl\{\lambda_{\min}(\Phi_j):\ j=\mathsf t(a)\ \text{for an arrow }a\text{ out of }i_0\bigr\};
\]
only targets with $W_j\neq0$ enter this minimum. If there are none,
then $s=0$, Lemma~\ref{lem:estimates}(iii) reads
$0\leq0$, and $\varpi_\star$, $\bar\varpi$ and the constants
\eqref{eq:constants} are not used.
Let $\Eone(b)=\int_1^\infty e^{-xb}\frac{dx}x$ be the exponential integral;
for $b>0$,
\begin{equation}\label{eq:E1}
 \frac{e^{-b}}{1+b}\leq\Eone(b)\leq e^{-b}\log\bigl(1+\tfrac1b\bigr),\qquad
 \Eone'(b)=-\frac{e^{-b}}b ,
\end{equation}
because the differences of the three functions tend to $0$ at infinity and
have derivatives of constant sign.

\begin{lemma}\label{lem:estimates}
For all charges and every metric:
\begin{enumerate}
\item[(i)] $\lambda$ is nondecreasing along the flow; its right derivative is
 the minimum of
 $\lambda\bigl(\inner v{\Phi_{AA}v}-\inner u{\Phi_Uu}\bigr)$ over the unit
 vectors $v$ with $F^\dagger Fv=\lambda v$, where $u=Fv/\sqrt\lambda$.
\item[(ii)] $\norm{E_R}\leq\sqrt{q_{\mathrm{out}}s/\lambda}$ and
 $\sum_j\norm{Y_j}_1\leq s$; hence, with $\mathsf d=\min(n_{i_0},r)$ and
 $\mathcal E_V:=\Omega_V-\diag(z_\ast+\mu_{i_0},\ z_\ast e^{-F^\dagger F})$,
 \[
 \norm{\mathcal E_V}_1\leq2\sqrt{\mathsf d\,q_{\mathrm{out}}s/\lambda}+r\kappa_0s .
 \]
\item[(iii)] $\dot s\leq\bigl(-\varpi_\star+\frac{4\pi q_{\mathrm{out}}}{y_\ast\lambda}\bigr)s$
 along the flow. Moreover $\lambda_{\min}(\Phi_j)\geq\arg(z_j+\cN)$ for
 each $j\neq i_0$ with $W_j\neq 0$. At every sink with $W_j\neq 0$ one has $\lambda_{\min}(\Phi_j)\geq\alpha_j$.
\item[(iv)] If $\sqrt{q_{\mathrm{out}}s/\lambda}\leq y_\ast/8$, then
 $\norm{\delta_{i_0}}\leq\frac{4\pi}{y_\ast}\sqrt{q_{\mathrm{out}}s/\lambda}$;
 and $\norm{\delta_j}\leq\norm{Y_j}/y_j\leq s/y_j$ for $j\neq i_0$.
\item[(v)] Along the flow
 \[
 \frac d{dt}\tr\Eone(L)\leq-\tr\bigl(e^{-L}|\Phi_U|\bigr),\qquad
 \tr\bigl(e^{-L}|\Phi_U|\bigr)\geq
 \frac{y_\ast(\tr e^{-L})^2}{(y_\ast+|x_\ast|)\tr e^{-L}+\tr J_+} .
 \]
\end{enumerate}
\end{lemma}

\begin{proof}
(i) By \eqref{eq:adapted},
$\frac d{dt}(F^\dagger F)=\frac12(\Phi_{AA}F^\dagger F+F^\dagger F\Phi_{AA})-F^\dagger\Phi_UF$.
For a unit vector $v$ with $F^\dagger Fv=\lambda v$ this gives
$\inner v{\frac d{dt}(F^\dagger F)v}=\lambda\inner v{\Phi_{AA}v}-\inner{Fv}{\Phi_UFv}$,
which is positive because $\Phi_{AA}>0>\Phi_U$; and the right derivative of
the least eigenvalue of a differentiable family of Hermitian matrices is the
minimum of these numbers over the unit vectors of the eigenspace.

(ii) Since $E_R=\sum_aH_a^\dagger X_a\,FR(F^\dagger F)$ and
$\norm{FR(F^\dagger F)}_{\mathrm{op}}\leq\lambda^{-1/2}$, the Cauchy--Schwarz
inequality gives the first bound. $\sum_j\tr Y_j=\tr\bigl((1-e^{-L})S\bigr)\leq s$.
The off-diagonal part of $\mathcal E_V$ has rank at most $2\mathsf d$ and
trace norm $2\norm{E_R}_1\leq2\sqrt{\mathsf d}\norm{E_R}$, and its remaining
block is $K_{F^\dagger F}(C_T)=K_{F^\dagger F}(F^\dagger SF)$; use
Proposition~\ref{prop:used}(v).

(iii) By \eqref{eq:adapted},
\[
 \dot s=\sum_a\tr(X_a^\dagger X_a\Phi_U)-\sum_a\tr(X_a^\dagger\Phi_{\mathsf t(a)}X_a)
 +2\Rea\sum_a\tr\bigl(X_a^\dagger H_a\Phi_{HA}F^{-1}\bigr).
\]
The first sum is nonpositive, the second is at least $\varpi_\star s$, and the last
term is at most $2\sqrt{s\,q_{\mathrm{out}}/\lambda}\,\norm{\Phi_{HA}}$. In
\eqref{eq:adaptedblocks} the block $z_\ast+\mu_{i_0}$ is normal with
imaginary part $y_\ast$, so its singular values are at least $y_\ast$, and
Proposition~\ref{prop:used}(iii) gives
$\norm{\Phi_{HA}}\leq\frac{2\pi}{y_\ast}\norm{E_R}\leq\frac{2\pi}{y_\ast}\sqrt{q_{\mathrm{out}}s/\lambda}$.
For the phases: $\Phi_j=g_j(M_j)$ with $g_j(u)=\arg(z_j+u)$ decreasing, and
$M_j\leq\sum_{a:j\to\cdot}\rho_a^\dagger\rho_a\leq\cN$ by \eqref{eq:blocks};
for a sink $M_j\leq0$.

(iv) The first statement is Proposition~\ref{prop:used}(iii) for
\eqref{eq:adaptedblocks}. For $j\neq i_0$, $\Phi_j=g_j(\mu_j-Y_j)$ and
$\arg(z_j+\mu_j)=g_j(\mu_j)$ by \eqref{eq:adaptedblocks}, $|g_j'|\leq1/y_j$, and a
function with Lipschitz constant $c$ satisfies
$\norm{g(A)-g(A')}\leq c\norm{A-A'}$ for Hermitian matrices $A,A'$ in the
Hilbert--Schmidt norm, as one sees by expanding in eigenbases of $A$ and
$A'$.

(v) By \eqref{eq:adapted},
$\dot L=F\Phi_{AA}F^\dagger-\frac12(\Phi_UL+L\Phi_U)$, so by \eqref{eq:E1}
\[
 \frac d{dt}\tr\Eone(L)=-\tr\bigl(e^{-F^\dagger F}\Phi_{AA}\bigr)+\tr\bigl(e^{-L}\Phi_U\bigr),
\]
and both terms are nonpositive. Let $(u_k)$ be an orthonormal eigenbasis of
$\Phi_U$ with eigenvalues $\psi_k\in(-\pi,0)$, and put
$e_k=\inner{u_k}{e^{-L}u_k}>0$, $j_k=\inner{u_k}{Ju_k}$. Since
$-\Omega_U=z_\ast e^{-L}-J$ has the polar phase $\Phi_U+\pi$ and a positive
definite modulus, $\inner{u_k}{(z_\ast e^{-L}-J)u_k}$ is a positive multiple
of $e^{\ii(\psi_k+\pi)}$; hence
$\cot(\psi_k+\pi)=(x_\ast e_k-j_k)/(y_\ast e_k)$ and
\[
 |\psi_k|=\operatorname{arccot}\frac{j_k-x_\ast e_k}{y_\ast e_k}
 \geq\operatorname{arccot}\frac{(j_k)_++|x_\ast|e_k}{y_\ast e_k}
 \geq\frac{y_\ast e_k}{(y_\ast+|x_\ast|)e_k+(j_k)_+},
\]
because $\operatorname{arccot}$ is decreasing and
$\operatorname{arccot}x\geq\frac1{1+x}$ for $x\geq0$. Therefore
\[
 \tr\bigl(e^{-L}|\Phi_U|\bigr)=\sum_k|\psi_k|e_k
 \geq\sum_k\frac{y_\ast e_k^2}{(y_\ast+|x_\ast|)e_k+(j_k)_+}
 \geq\frac{y_\ast(\sum_ke_k)^2}{(y_\ast+|x_\ast|)\sum_ke_k+\sum_k(j_k)_+}
\]
by the Cauchy--Schwarz inequality, and $\sum_ke_k=\tr e^{-L}$,
$\sum_k(j_k)_+\leq\tr J_+$.
\end{proof}

\subsection{Constants}
Fix a height function $w$ and let $\gamma_w$, $c_Q$, $\bar\cN$, $t_0$ be as in
Theorem~\ref{thm:energy}. We assume $c_Q>0$, that is, $W_j\neq0$ for some
$j\neq i_0$; otherwise all arrows vanish, Proposition~\ref{prop:sink} below
applies, and we put $\bar\cN=0$ and $t_0=0$. Let $\bar\varpi$ be the least of the
numbers
\[
 \arg(z_j+\bar\cN)\ \ (j\text{ not a sink}),\qquad\alpha_j\ \ (j\text{ a sink}),
\]
over the nonzero target spaces $W_j$ of arrows out of $i_0$; by
Lemma~\ref{lem:estimates}(iii) and Theorem~\ref{thm:energy},
$\varpi_\star\geq\bar\varpi$ for $t\geq t_0$. Let $\lambda_0=\lambda(0)$,
$c_\lambda=1/\log(1+1/\lambda_0)$, $\Xi_0=\tr\Eone(L)|_{t=0}$, and
\begin{equation}\label{eq:constants}
 \begin{gathered}
 \kappa_1=\frac{y_\ast\,c_\lambda^2}{\sqrt2\,|z_\ast|\,c_\lambda\,\Xi_0+\frac75\bar\cN},\qquad
 \Lambda_1=\max\Bigl\{1,\frac{8\pi\bar\cN}{\bar\varpi\,y_\ast}\Bigr\},\\
 t_1=t_0+\frac{(1+\Lambda_1)e^{\Lambda_1}}{\kappa_1},\qquad
 s_1=\frac{\bar\cN}{1-e^{-1}},\qquad
 \mathfrak c_1=2\sqrt{\mathsf d\,\bar\cN s_1}+r\kappa_0s_1,\\
 \mathfrak c_2=\frac\pi{\sqrt{|z_\ast|}}\bigl(\sqrt{\mathfrak c_1}+\sqrt{rs_1}\bigr),\qquad
 \mathfrak c_3=\max\Bigl\{0,2\log\frac{2\mathfrak c_2}\pi\Bigr\},\\
 C_1=\mathfrak c_1+(r+1)s_1+2|z_\ast|re^{\mathfrak c_3},\qquad
 t_2=t_1+\frac2{\bar\varpi}\log^+\frac{8\bar\varpi\,s_1}{\pi y_\ast}.
 \end{gathered}
\end{equation}
In these constants the subscripts are ordinal labels and do not refer to
vertices.

\begin{theorem}[Decay of the exact part]\label{thm:decay}
Let $\mathsf E$ be a complex with a surjective differential at $i_0$ and
$r\geq1$, with at least one arrow from $i_0$ to a vertex $j$ with
$W_j\neq0$; fix a height function $w$ and let the constants be as in
\eqref{eq:constants}. For all
charges in the upper half-plane and along every trajectory of the phase flow:
\begin{enumerate}
\item[(i)] $\lambda(t)\to\infty$; for $t\geq t_1$, $\lambda(t)\geq
 \max\{\Lambda_1,\ \frac{\bar\varpi}4(t-t_1)-\mathfrak c_3\}$ and
 $s(t)\leq s_1e^{-\frac{\bar\varpi}2(t-t_1)}$;
\item[(ii)] for $t\geq t_1$,
 \[
 \norm{\Omega_V-(z_\ast+\mu_{i_0})\oplus0}_1+\norm{\Omega_U}_1
 +\sum_{j\neq i_0}\norm{\Omega_j-(z_j+\mu_j)}_1
 \leq C_1\,e^{-\frac{\bar\varpi}4(t-t_1)};
 \]
 the difference between the mass and the mass $\sum_i\norm{z_i+\mu_i}_1$ of
 the cohomology representation with its induced metrics satisfies the same
 bound;
\item[(iii)] for $t\geq t_2$,
 $\norm{\delta_{i_0}}^2\leq\frac{2\pi\bar\varpi}{y_\ast}\,s$ and
 $\norm{\delta_j}\leq s/y_j$ for $j\neq i_0$.
\end{enumerate}
In particular the exact part decays.
\end{theorem}

The theorem makes no assumption on the cohomology representation.

\begin{proof}
\emph{Growth of $\lambda$.} Put $\Xi=\tr\Eone(L)$. By
Lemma~\ref{lem:estimates}(i) all eigenvalues of $L$ are at least $\lambda_0$
for $t\geq0$, so \eqref{eq:E1} gives $\tr e^{-L}\geq c_\lambda\Xi$. For
$t\geq t_0$, $\tr J_+=\tr J+\tr J_-\leq\frac75\tr J\leq\frac75\bar\cN$ by
Proposition~\ref{prop:used}(i), \eqref{eq:energybounds} and
Theorem~\ref{thm:energy}, and $y_\ast+|x_\ast|\leq\sqrt2|z_\ast|$. The
function $x\mapsto x^2/(cx+d)$ is increasing for $x>0$, and $\Xi$ is
nonincreasing; so Lemma~\ref{lem:estimates}(v) gives, for $t\geq t_0$,
\[
 \dot\Xi\leq-\frac{y_\ast c_\lambda^2\,\Xi^2}{\sqrt2|z_\ast|c_\lambda\,\Xi+\frac75\bar\cN}
 \leq-\kappa_1\Xi^2,
\]
hence $\Eone(\lambda(t))\leq\Xi(t)<1/(\kappa_1(t-t_0))$ for $t>t_0$. For
$t\geq t_1$ this is at most $e^{-\Lambda_1}/(1+\Lambda_1)\leq\Eone(\Lambda_1)$
by \eqref{eq:E1}, so $\lambda(t)\geq\Lambda_1$.

\emph{Decay of $s$.} For $t\geq t_0$ we have $q_{\mathrm{out}}\leq\bar\cN$ and
$\varpi_\star\geq\bar\varpi$, so for $t\geq t_1$,
$\frac{4\pi q_{\mathrm{out}}}{y_\ast\lambda}\leq\frac{\bar\varpi}2$ and
$\dot s\leq-\frac{\bar\varpi}2s$ by Lemma~\ref{lem:estimates}(iii); at $t_1$,
$\lambda\geq1$ and $s\leq\bar\cN/(1-e^{-\lambda})\leq s_1$ by
\eqref{eq:energybounds}.

\emph{Linear growth of $\lambda$.} By Lemma~\ref{lem:estimates}(ii) and
Proposition~\ref{prop:used}(v), for $t\geq t_1$,
$\norm{\mathcal E_V}_1\leq\mathfrak c_1e^{-\frac{\bar\varpi}4(t-t_1)}$ and
$\norm J_1\leq rs_1e^{-\frac{\bar\varpi}2(t-t_1)}$. Let $v$ be a unit vector of
$(\ker D)^\perp$ with $F^\dagger Fv=\lambda v$ and $u=Fv/\sqrt\lambda$. Apply
Proposition~\ref{prop:used}(iv) to $T=\Omega_V$ with
$T_\circ=\diag(z_\ast+\mu_{i_0},\ z_\ast e^{-F^\dagger F})$, whose first
block is normal and invertible and for which
$|T_\circ|v=|z_\ast|e^{-\lambda}v$ and $\Psi_\circ v=\alpha_\ast v$, and to
$T=-\Omega_U=z_\ast e^{-L}-J$ with $T_\circ=z_\ast e^{-L}$ and the vector
$u$:
\[
 |\inner v{(\Phi_V-\alpha_\ast)v}|\leq\pi\sqrt{\frac{\norm{\mathcal E_V}_1e^\lambda}{|z_\ast|}},
 \qquad
 |\inner u{(\Phi_U+\pi-\alpha_\ast)u}|\leq\pi\sqrt{\frac{\norm J_1e^\lambda}{|z_\ast|}} .
\]
By Lemma~\ref{lem:estimates}(i) the right derivative of $\lambda$ satisfies
$\dot\lambda\geq\lambda\bigl[\pi-\mathfrak c_2e^{\lambda/2-\frac{\bar\varpi}8(t-t_1)}\bigr]$
for $t\geq t_1$. Where $\lambda$ is at most
$\frac{\bar\varpi}4(t-t_1)-\mathfrak c_3$, the bracket is at least $\frac\pi2$, so
$\dot\lambda\geq\frac\pi2\lambda\geq\frac\pi2>\frac{\bar\varpi}4$; since
$\lambda(t_1)\geq1>-\mathfrak c_3$, the graph of $\lambda$ cannot cross below
this line. This proves (i).

(ii) By \eqref{eq:adaptedblocks},
$\norm{\Omega_V-(z_\ast+\mu_{i_0})\oplus0}_1\leq\norm{\mathcal E_V}_1+|z_\ast|re^{-\lambda}$,
$\norm{\Omega_U}_1\leq rs+|z_\ast|re^{-\lambda}$ and
$\sum_j\norm{\Omega_j-(z_j+\mu_j)}_1=\sum_j\tr Y_j\leq s$; by (i),
$e^{-\lambda}\leq e^{\mathfrak c_3}e^{-\frac{\bar\varpi}4(t-t_1)}$. The statement
on the masses follows from the triangle inequality for the trace norm.

(iii) For $t\geq t_2$, $s\leq\frac{\pi y_\ast}{8\bar\varpi}$ by (i). Then
$\frac{q_{\mathrm{out}}s}{\lambda y_\ast^2}\leq\frac{\bar\cN}{\Lambda_1}\frac s{y_\ast^2}
\leq\frac{\bar\varpi\,s}{8\pi y_\ast}\leq\frac1{64}$, and
Lemma~\ref{lem:estimates}(iv) gives
$\norm{\delta_{i_0}}^2\leq\frac{16\pi^2}{y_\ast^2}\frac{q_{\mathrm{out}}s}\lambda
\leq\frac{2\pi\bar\varpi}{y_\ast}s$.
\end{proof}

\begin{proposition}[No arrow out of $i_0$]\label{prop:sink}
Let $\mathsf E$ have a surjective differential at $i_0$ and $r\geq1$, and
let $\rho_a=0$ for every arrow $a$ out of $i_0$; this holds if no arrow leaves
$i_0$, and if $W_j=0$ for all $j\neq i_0$. Then, for every metric, the complex is the orthogonal
direct sum of its cohomology representation and of the complex
$(\ker D)^\perp\to U$ at $i_0$, the flow preserves this decomposition,
$\delta_i=0$ for all $i$, the phases on $(\ker D)^\perp$ and $U$ are
$\alpha_\ast$ and $\alpha_\ast-\pi$, and
\[
 F(t)=e^{\pi t/2}F(0),\qquad
 \Omega_V=(z_\ast+\mu_{i_0})\oplus z_\ast e^{-F^\dagger F},\qquad
 \Omega_U=-z_\ast e^{-L}
\]
for all $t$. In particular
$\norm{\Omega_V-(z_\ast+\mu_{i_0})\oplus0}_1+\norm{\Omega_U}_1
\leq2|z_\ast|r\exp(-\lambda_0e^{\pi t})$, and the exact part decays. This
bound is at most $2|z_\ast|r\,e^{-(t-t_1)}$ for
$t\geq t_1:=\frac1\pi\log^+\frac1{\lambda_0}$.
\end{proposition}

\begin{proof}
Here $C_{\mathrm{out}}=0$, so $J=0$, $E_R=0$, $C_T=0$ and $Y_j=0$ in
\eqref{eq:blocks} and \eqref{eq:adaptedblocks}. The blocks of $\Omega_V$ are
$z_\ast+\mu_{i_0}$ and $z_\ast e^{-F^\dagger F}$, so $\Phi_{HA}=0$,
$\Phi_{AA}=\alpha_\ast$, $\Phi_U=\alpha_\ast-\pi$, and \eqref{eq:adapted} gives
$\dot F=\frac\pi2F$.
\end{proof}

\begin{remark}\label{rem:decay}
(1) The time $t_1$ contains the factor $e^{\Lambda_1}$, and
$\Lambda_1\geq8\bar\cN/y_\ast$.

(2) When an arrow leaves $i_0$, the phases $\Phi_{AA}$ and $\Phi_U$ on the
exact part need not tend to $\alpha_\ast$ and $\alpha_\ast-\pi$, in contrast
with Proposition~\ref{prop:sink}. In the proof of Theorem~\ref{thm:decay}(i), the bound
$\inner v{\Phi_Vv}-\inner u{\Phi_Uu}\geq\pi/2$ is used at a possible
contact with the line $\lambda=\frac{\bar\varpi}4(t-t_1)-\mathfrak c_3$.
It gives the stated lower bound for $\lambda(t)$.
\end{remark}
